% Source: https://artofproblemsolving.com/community/c3h1715918_junior_regional_olympiad__fbh_2015_grade_8_problem_5
% Stable identity: aops-topic-1715918-q01
Colour the elements of one set red and of the other blue, and suppose, for a contradiction, that there are no distinct $x<y$ of the same colour with $x+y$ also of that colour (call such $x,y,x+y$ a \emph{monochromatic sum}). Swapping the colours if necessary, $1$ is red.

\textbf{Case 1: $2$ is red.} Then $3=1+2$ is blue.
\begin{itemize}
\item If $4$ is red, then $5=1+4$ and $6=2+4$ are blue; then $8=3+5$ and $9=3+6$ are red, and $1+8=9$ is a red monochromatic sum. Contradiction.
\item If $4$ is blue, then $7=3+4$ is red, so $8=1+7$ and $9=2+7$ are blue. Now $5$ cannot be blue ($4+5=9$) and cannot be red ($2+5=7$). Contradiction.
\end{itemize}

\textbf{Case 2: $2$ is blue.}
\begin{itemize}
\item If $3$ is red, then $4=1+3$ is blue, so $6=2+4$ is red; then $7=1+6$ and $9=3+6$ are blue, and $2+7=9$ is a blue monochromatic sum. Contradiction.
\item If $3$ is blue, then $5=2+3$ is red, so $4$ is blue (as $1+4=5$) and $6$ is blue (as $1+5=6$). Then $2+4=6$ is a blue monochromatic sum. Contradiction.
\end{itemize}
Every case leads to a contradiction, so one of the two sets contains distinct $x,y$ and $x+y$. $\blacksquare$

(The number $9$ cannot be lowered: $\{1,2,4,8\}$ and $\{3,5,6,7\}$ split $\{1,\dots,8\}$ with no such triple.)
